大语言模型（LLM）驱动的命令行编码代理在长程软件开发任务中普遍面临一个被忽视的工程难题：代理会频繁停下来等待人工介入——等待“继续”、等待权限确认、等待菜单选择，或卡死在瞬时错误上。本文介绍 WhatyTerm，一个为消除这些中断点而设计的 AI 终端会话编排系统，并\textbf{以真实运行数据为核心}报告其设计意图与实际行为之间的落差。系统设计上包含三项机制：以规则层优先、LLM 兜底的分层决策；联合 Hook 事件、终端抓屏与进程指纹的三通道感知；以及 micro/meso/macro 三层嵌套的自主交付循环。然而，对一次真实部署（跨三天、四个会话、46 次自动决策）的日志分析揭示了一个与设计初衷相悖的现象：\textbf{全部 46 次决策 100\% 由零 Token 的规则层做出，LLM 兜底层一次也未被触发；其中 87\% 的动作是机械地发送“继续”，仅 1 次注入了真实的任务指令。}换言之，一个本应智能编排的自主开发系统，在真实运行中退化成了一个“自动点继续”的宏。我们剖析了这一退化的代码级成因——规则层在检测到空闲提示符时无条件发送“继续”，缺乏“本次继续是否推进了任务”的判断，也没有将该判断移交 LLM 的路径——并据此讨论自主代理 harness 中“规则层与推理层边界”这一更普遍的设计陷阱。本文的立场是：报告一个系统\emph{未按预期工作}的真实证据与根因，比宣称其设计优越更有价值。

\vspace{4pt}
\noindent\textbf{关键词：} AI 编码代理；自主软件开发；终端自动化；负面结果；人在环
